%% PACKAGES
% Geometry and Paper Setup
\documentclass{article}
\usepackage[a4paper, total={6in,8in}]{geometry}
\geometry{left=2cm,right=2cm,top=3cm,bottom=3cm}
\usepackage{graphicx}

% Mathematical Formatting
\usepackage{amsmath}
\usepackage{amsfonts}
\usepackage{amsthm}
\usepackage{amssymb}
\usepackage{blkarray}
\usepackage{multirow}
\usepackage{array}
\usepackage[most]{tcolorbox}
\usepackage{tikz}
\usepackage{tikz-cd}
\usetikzlibrary{arrows.meta, automata, positioning, decorations.pathreplacing}
\usepackage{forest}
\usepackage{dsfont}
\newcommand{\llb}{[\![}
\newcommand{\rrb}{]\!]}

% Code Listing and Formatting
\usepackage{listings}
\usepackage{color}
\usepackage{placeins}

% Text and Formatting
\usepackage{enumitem}
\usepackage[none]{hyphenat} % May cause issues with vspace in itemizemacro
\usepackage{subcaption}
\usepackage{xparse}
\usepackage{cancel}

%% THEOREMS AND ENVIRONMENTS
\theoremstyle{definition}
\newtheorem{definition}{Definition}[section]
\newtheorem{algorithm}{Algorithm}
\newtheorem{hypothesis}{Hypothesis}
\newtheorem{paradox}{Paradox}
\newtheorem{construction}{Construction}
\newtheorem{type}{Type}
\newtheorem{axiom}{Axiom}
\newtheorem*{answer}{Answer}
\newtheorem*{explanation}{Explanation}
\newtheorem{criterion}{Criterion}
\newtheorem{proposition}[definition]{Proposition}
\newtheorem{theorem}[definition]{Theorem}
\newtheorem{corollary}[definition]{Corollary}
\newtheorem{lemma}[definition]{Lemma}
\newtheorem{conjecture}[definition]{Conjecture}
\theoremstyle{remark}
\newtheorem{example}[definition]{Example}
\newtheorem*{example2}{Example}
\theoremstyle{plain}

\newtcolorbox[number within=section]{theorembox}[1][]{
	breakable,
	enhanced,
	arc=0mm,
	colback=white,
	colframe=white, % Remove frame color
	coltitle=black,
	fonttitle=\bfseries,
	title={\refstepcounter{definition}Theorem \thetheorem. #1}, % Increment theorem counter
	attach title to upper,
	before upper={\tcbtitle\quad},
	after upper={\hfill}, % Remove the arrow
	sharp corners,
	boxsep=2pt, % Add some padding
	left=10pt,   % Add space after the left black bar
	borderline west={2pt}{0pt}{black}, % Add the left border
}

\newtcolorbox[number within=section]{definitionbox}[1][]{
	breakable,
	enhanced,
	arc=0mm,
	colback=white,
	colframe=white, % Remove frame color
	coltitle=black,
	fonttitle=\bfseries,
	title={\refstepcounter{definition}Definition \thedefinition. #1}, % Increment theorem counter
	attach title to upper,
	before upper={\tcbtitle\quad},
	after upper={\hfill}, % Remove the arrow
	sharp corners,
	boxsep=2pt, % Add some padding
	left=10pt,   % Add space after the left black bar
	borderline west={2pt}{0pt}{black}, % Add the left border
}

\newtcolorbox[number within=section]{propositionbox}[1][]{
	breakable,
	enhanced,
	arc=0mm,
	colback=white,
	colframe=white, % Remove frame color
	coltitle=black,
	fonttitle=\bfseries,
	title={\refstepcounter{definition}Proposition \theproposition. #1}, % Increment theorem counter
	attach title to upper,
	before upper={\tcbtitle\quad},
	after upper={\hfill}, % Remove the arrow
	sharp corners,
	boxsep=2pt, % Add some padding
	left=10pt,   % Add space after the left black bar
	borderline west={2pt}{0pt}{black}, % Add the left border
}

\newtcolorbox[number within=section]{corollarybox}[1][]{
	breakable,
	enhanced,
	arc=0mm,
	colback=white,
	colframe=white, % Remove frame color
	coltitle=black,
	fonttitle=\bfseries,
	title={\refstepcounter{definition}Corollary \thecorollary. #1}, % Increment theorem counter
	attach title to upper,
	before upper={\tcbtitle\quad},
	after upper={\hfill}, % Remove the arrow
	sharp corners,
	boxsep=2pt, % Add some padding
	left=10pt,   % Add space after the left black bar
	borderline west={2pt}{0pt}{black}, % Add the left border
}

\newtcolorbox[number within=section]{lemmabox}[1][]{
	breakable,
	enhanced,
	arc=0mm,
	colback=white,
	colframe=white, % Remove frame color
	coltitle=black,
	fonttitle=\bfseries,
	title={\refstepcounter{definition}Lemma \thelemma. #1}, % Increment theorem counter
	attach title to upper,
	before upper={\tcbtitle\quad},
	after upper={\hfill}, % Remove the arrow
	sharp corners,
	boxsep=2pt, % Add some padding
	left=10pt,   % Add space after the left black bar
	borderline west={2pt}{0pt}{black}, % Add the left border
}

\newtcolorbox[number within=section]{conjecturebox}[1][]{
	breakable,
	enhanced,
	arc=0mm,
	colback=white,
	colframe=white, % Remove frame color
	coltitle=black,
	fonttitle=\bfseries,
	title={\refstepcounter{definition}Conjecture \theconjecture. #1}, % Increment theorem counter
	attach title to upper,
	before upper={\tcbtitle\quad},
	after upper={\hfill}, % Remove the arrow
	sharp corners,
	boxsep=2pt, % Add some padding
	left=10pt,   % Add space after the left black bar
	borderline west={2pt}{0pt}{black}, % Add the left border
}

\newtcolorbox[number within=section]{questionbox}[1][]{
	breakable,
	enhanced,
	arc=0mm,
	colback=white,
	colframe=white, % Remove frame color
	coltitle=black,
	fonttitle=\bfseries,
	title={Question.}, % Increment theorem counter
	attach title to upper,
	before upper={\tcbtitle\quad},
	after upper={\hfill}, % Remove the arrow
	sharp corners,
	boxsep=2pt, % Add some padding
	left=10pt,   % Add space after the left black bar
	borderline west={2pt}{0pt}{black}, % Add the left border
}

\newtcolorbox{algorithmbox}[1][]{
	breakable,
	enhanced,
	arc=0mm,
	colback=white,
	colframe=white, % Remove frame color
	coltitle=black,
	fonttitle=\bfseries,
	title={\refstepcounter{algorithm}Algorithm \thealgorithm. #1}, % Increment theorem counter
	attach title to upper,
	before upper={\tcbtitle\quad},
	after upper={\hfill}, % Remove the arrow
	sharp corners,
	boxsep=2pt, % Add some padding
	left=10pt,   % Add space after the left black bar
	borderline west={2pt}{0pt}{black}, % Add the left border
}

\newtcolorbox[number within=section]{leftbar}[1][]{
	breakable,
	enhanced,
	arc=0mm,
	colback=white,
	colframe=white, % Remove frame color
	coltitle=black,
	fonttitle=\bfseries,
	sharp corners,
	boxsep=2pt, % Add some padding
	left=10pt,   % Add space after the left black bar
	borderline west={2pt}{0pt}{black}, % Add the left border
}

\newtcolorbox[number within=section]{boxedup}[1][]{
	breakable,
	enhanced,
	arc=0mm,
	colback=white,
	colframe=black, % Set frame color to black
	coltitle=black,
	fonttitle=\bfseries,
	sharp corners,
	boxsep=2pt, % Add some padding
	borderline={0.5pt}{0pt}{black}, % Add the complete border around the box
}

%% MATH OPERATORS
\DeclareMathOperator{\Span}{span}
\DeclareMathOperator{\comp}{comp}
\DeclareMathOperator{\proj}{proj}
\DeclareMathOperator{\Tr}{Tr}
\DeclareMathOperator{\size}{size}
\DeclareMathOperator{\dom}{dom}
\DeclareMathOperator{\range}{range}
\DeclareMathOperator{\rank}{rank}
\DeclareMathOperator{\nullity}{nullity}
\DeclareMathOperator{\id}{id}
\DeclareMathOperator{\im}{im}
\DeclareMathOperator{\lcm}{lcm}
\DeclareMathOperator{\ord}{ord}
\DeclareMathOperator{\Null}{Null}
\DeclareMathOperator{\argmin}{argmin}
\DeclareMathOperator{\argmax}{argmax}
\DeclareMathOperator{\Extr}{Extr}
\DeclareMathOperator{\extr}{extr}
\DeclareMathOperator{\Maps}{Maps}
\DeclareMathOperator{\Perm}{Perm}
\DeclareMathOperator{\Aut}{Aut}
\DeclareMathOperator{\Inn}{Inn}
\DeclareMathOperator{\Out}{Out}
\DeclareMathOperator{\Syl}{Syl}
\DeclareMathOperator{\Char}{char}
\DeclareMathOperator{\Dih}{Dih}
\DeclareMathOperator{\Dic}{Dic}
\DeclareMathOperator{\Mat}{Mat}
\DeclareMathOperator{\cost}{cost}
\DeclareMathOperator{\Var}{Var}
\DeclareMathOperator{\Cov}{Cov}
\DeclareMathOperator{\Bernoulli}{Bernoulli}
\DeclareMathOperator{\Geo}{Geo}
\DeclareMathOperator{\Po}{Po}
\DeclareMathOperator{\Exp}{Exp}
\DeclareMathOperator{\Lognormal}{Lognormal}
\DeclareMathOperator{\conv}{conv}
\DeclareMathOperator{\sgn}{sgn}
\DeclareMathOperator{\ReP}{Re}
\DeclareMathOperator{\ImP}{Im}
\DeclareMathOperator{\Arg}{Arg}
\DeclareMathOperator{\Log}{Log}
\DeclareMathOperator{\pr}{pr}
\DeclareMathOperator{\sech}{sech}
\DeclareMathOperator{\csch}{csch}
\DeclareMathOperator{\arcsinh}{arcsinh}
\DeclareMathOperator{\arccosh}{arccosh}
\DeclareMathOperator{\arctanh}{arctanh}
\DeclareMathOperator{\Int}{int}
\DeclareMathOperator{\Res}{Res}
\DeclareMathOperator{\PV}{PV}
\DeclareMathOperator{\Free}{Free}
\DeclareMathOperator{\Th}{Th}
\DeclareMathOperator{\Cn}{Cn}
\DeclareMathOperator{\Mod}{Mod}
\DeclareMathOperator{\esssup}{ess \, sup}
\DeclareMathOperator{\Lip}{Lip}
\DeclareMathOperator{\Orb}{Orb}
\DeclareMathOperator{\Stab}{Stab}
\DeclareMathOperator{\Fix}{Fix}
\DeclareMathOperator{\bdim}{bdim}
\DeclareMathOperator{\SBM}{SBM}
\DeclareMathOperator{\SSBM}{SSBM}
\DeclareMathOperator{\HSBM}{HSBM}
\DeclareMathOperator{\diam}{diam}
\DeclareMathOperator{\rad}{rad}
\newcommand{\herm}{\mathsf{H}}
\newcommand{\Dashv}{\mathrel{\reflectbox{$\vDash$}}}

%% COMMANDS AND SETTINGS
\newcommand\hcancel[2][black]{\setbox0=\hbox{$#2$}
\rlap{\raisebox{.45\ht0}{\textcolor{#1}{\rule{\wd0}{1pt}}}}#2}

\newcommand{\ind}{\mathrel{\perp\!\!\!\perp}}

\newcommand\twoheaduparrow{\mathrel{\rotatebox{0}{$\twoheadrightarrow$}}}

\definecolor{dkgreen}{rgb}{0,0.6,0}
\definecolor{gray}{rgb}{0.5,0.5,0.5}
\definecolor{mauve}{rgb}{0.58,0,0.82}

\lstset{frame=tb,
	language=Python,
	aboveskip=3mm,
	belowskip=3mm,
	showstringspaces=false,
	columns=flexible,
	basicstyle={\small\ttfamily},
	numbers=none,
	numberstyle=\tiny\color{gray},
	keywordstyle=\color{blue},
	commentstyle=\color{dkgreen},
	stringstyle=\color{mauve},
	breaklines=true,
	breakatwhitespace=true,
	tabsize=3,
	literate={-}{{-}}1
	{“}{{``}}1
	{”}{{''}}1
	{–}{{-\ }}1
	{—}{{---}}1
	{---}{{---\ }}3
    {&}{{\&}}1
    {$}{{\$}}1
}

% For all-black codeblock
%\lstset{frame=tb,
%	language=Python,
%	aboveskip=3mm,
%	belowskip=3mm,
%	showstringspaces=false,
%	columns=flexible,
%	basicstyle={\small\ttfamily},
%	numbers=none,
%	numberstyle=\tiny\color{black},
%	keywordstyle=\color{black},
%	commentstyle=\color{black},
%	stringstyle=\color{black},
%	breaklines=true,
%	breakatwhitespace=true,
%	tabsize=3,
%	literate={-}{{-}}1
%	{“}{{``}}1
%	{”}{{''}}1
%	{–}{{-\ }}1
%	{—}{{---}}1
%	{---}{{---\ }}3
%    {&}{{\&}}1
%    {$}{{\$}}1
%}

\setcounter{tocdepth}{2}
\setlength{\parindent}{0pt}
\renewcommand{\labelitemi}{\tiny\textbullet}

\setcounter{topnumber}{10}    % Maximum floats at the top of a page
\setcounter{bottomnumber}{10} % Maximum floats at the bottom of a page
\setcounter{totalnumber}{20}  % Maximum floats on a page
\renewcommand{\topfraction}{0.95}   % Allow up to 95% of a page for floats at the top
\renewcommand{\bottomfraction}{0.95} % Allow up to 95% of a page for floats at the bottom
\renewcommand{\textfraction}{0.05}  % Allow as little as 5% text on a page
\renewcommand{\floatpagefraction}{0.9} % Require 90% float content to create a float-only page

%% ENVIRONMENTS
\newenvironment{remark}
{\par\addvspace{\topsep}\leftskip=2em\textit{Remark.}\space}
{\par\addvspace{\topsep}}

% vspace may be needed because of hyphenat
\newenvironment{enumeratemacro}{
    \vspace{-1ex}
    \begin{enumerate}[labelsep=0.5em, leftmargin=2em, itemsep=-0.5ex]
}{
    \end{enumerate}
}

\newenvironment{itemizemacro}{
    \vspace{-1ex}
    \begin{itemize}[label=\textbullet, labelsep=0.5em, leftmargin=2em, itemsep=-0.5ex]
}{
    \end{itemize}
}

\newenvironment{itemizesublistmacro}{
    \vspace{-1.5ex}
    \begin{itemize}[label=\textendash, labelsep=0.5em, leftmargin=2em, itemsep=-0.5ex]
}{
    \end{itemize}
    \vspace{-1ex}
}

\NewDocumentCommand{\figuremacro}{O{250px} m}
{
    \begin{figure}[ht]
        \centering
        \includegraphics[width=#1]{#2}
    \end{figure}
    \FloatBarrier
}